% =====================================================================
%  Préambule commun — CV Rachidi Adébissi IDONIYI
%  Design partagé par toutes les versions (FullStack, DataScience, EN).
%  Ne contient AUCUN contenu : uniquement la mise en page et les macros.
% =====================================================================
\documentclass[11pt,a4paper]{article}
\makeatletter

% ---------- Compatibilité ATS : texte sélectionnable et copiable ----------
\pdfgentounicode=1
\input{glyphtounicode}

\usepackage[T1]{fontenc}
\usepackage[utf8]{inputenc}
\usepackage{lmodern}                 % Computer Modern (même rendu que le CV d'origine)
\usepackage[babel=true]{microtype}   % meilleure justification, moins de débordements
\ifdefined\cvlangen
  \usepackage[english]{babel}
\else
  \usepackage[french]{babel}
\fi

\usepackage[a4paper,left=1.8cm,right=1.8cm,top=1.3cm,bottom=1.25cm]{geometry}
\usepackage{xcolor}
\usepackage{fontawesome5}
\usepackage{accsupp}                 % icônes ignorées par les ATS (ActualText vide)
\usepackage{etoolbox}
\usepackage{enumitem}
\usepackage{titlesec}
\usepackage{tabularx}
\usepackage{needspace}
\usepackage{qrcode}
\usepackage[hidelinks,unicode]{hyperref}
\urlstyle{same}

% ---------- Couleurs (accent du CV d'origine conservé) ----------
\definecolor{accent}{HTML}{6C63FF}      % couleur d'origine : titres, icônes, filets
\definecolor{accentfonce}{HTML}{4A42D6} % même teinte, plus foncée : liens (lisibles imprimés)
\definecolor{encre}{HTML}{1E1E28}       % texte principal
\definecolor{sourd}{HTML}{5A5A66}       % texte secondaire (dates, établissements)
\definecolor{todo}{HTML}{FFE45C}        % surlignage des données manquantes (brouillon)

\color{encre}
\setlength{\parindent}{0pt}
\setlength{\parskip}{0pt}
\linespread{1.05}
\pagestyle{empty}
\raggedbottom

% =====================================================================
%  Données à compléter : \cvset{clé}{valeur} dans donnees.tex
%  - Version FINALE   : une phrase dont la donnée est vide est MASQUÉE.
%  - Version BROUILLON: la donnée manquante apparaît surlignée en jaune.
%  Le mode brouillon s'active en compilant avec \def\cvdraftmode{} (voir build.sh).
% =====================================================================
\newif\ifcvdraft
\ifdefined\cvdraftmode \cvdrafttrue \fi

\newcommand\cvset[2]{\csdef{cvd@#1}{#2}}
\newcommand\cvget[1]{\csuse{cvd@#1}}
\newcommand\cvtodobox[1]{\ifcvdraft{\setlength\fboxsep{1pt}\colorbox{todo}{\color{black}\sffamily\footnotesize\strut[#1]}}\fi}

% \TODO{description}           -> surligné en brouillon, rien en final
% \TODO[clé]{description}      -> valeur de la clé si renseignée, sinon comme ci-dessus
\NewDocumentCommand\TODO{o m}{%
  \IfNoValueTF{#1}{\cvtodobox{#2}}{\ifcsvoid{cvd@#1}{\cvtodobox{#2}}{\cvget{#1}}}}

% \cvif{clé}{texte}    -> texte affiché si la clé est renseignée (ou en brouillon)
\newcommand\cvif[2]{\ifcsvoid{cvd@#1}{\ifcvdraft#2\fi}{#2}}
% \cvifoui{clé}{texte} -> texte affiché si la clé vaut « oui » (ou en brouillon)
\newcommand\cvifoui[2]{\ifcsstring{cvd@#1}{oui}{#2}{\ifcvdraft#2\fi}}

% ---------- Icônes (invisibles pour les ATS) ----------
\newcommand\cvicon[1]{\BeginAccSupp{method=escape,ActualText={}}{\color{accent}#1}\EndAccSupp{}}

% ---------- Liens : texte visible + cliquable ----------
\newcommand\cvlink[2]{\href{#1}{\mbox{\color{accentfonce}#2}}}
\newcommand\cvmail[1]{\edef\cv@tmp{\noexpand\href{mailto:#1}}\cv@tmp{{\color{accentfonce}#1}}}
% Ligne de lien sous un projet : \cvweb{libellé}{url}{texte affiché}
\newcommand\cvweb[3]{{\small\cvicon{\faExternalLink*}~#1~\cvlink{#2}{#3}}\par}

% ---------- Titres de section : icône + titre coloré + filet ----------
\titleformat{\section}{\Large\bfseries\color{accent}}{}{0pt}{}
  [\vspace{2pt}{\color{accent!70}\titlerule[0.6pt]}]
\titlespacing*{\section}{0pt}{10pt}{5.5pt}
\newcommand\cvsection[2]{\section*{\cvicon{#1}\hspace{7pt}#2}}

% ---------- Listes (tiret d'origine, en couleur d'accent) ----------
\setlist[itemize]{label={\color{accent}\textendash}, leftmargin=1.15em, labelsep=0.5em,
  itemsep=1.6pt, parsep=0pt, topsep=4pt, partopsep=0pt}

% Bloc insécable : une entrée et sa liste ne sont jamais coupées entre deux pages
\newenvironment{cvblock}{\par\noindent\begin{minipage}[t]{\linewidth}}{\end{minipage}\par\vspace{7pt}}

% \cventry{intitulé}{dates}{organisation}{lieu}  — format d'origine sur deux lignes
\newcommand\cventry[4]{%
  {\textbf{#1}}\hfill{\itshape\color{sourd}#2}\par
  \ifblank{#3#4}{}{{\itshape\color{sourd}#3}\hfill{\itshape\color{sourd}#4}\par}}

% \cvproject{nom}{technologies}
\newcommand\cvproject[2]{{\textbf{#1}}\hfill{\itshape\color{sourd}#2}\par\vspace{1.5pt}}
\newcommand\cvdesc[1]{#1\par}

% Lignes de compétences
\newenvironment{cvskills}{\renewcommand\arraystretch{1.32}\tabularx{\linewidth}{@{}>{\bfseries}p{4.15cm}@{}X@{}}}{\endtabularx\par}

% ---------- En-tête avec QR code vers le portfolio ----------
% \cventete{nom}{titre}{ligne de contacts 1}{ligne de contacts 2}{disponibilité}{légende QR}
\newcommand\cventete[6]{%
  \noindent
  \begin{minipage}[c]{2.3cm}\mbox{}\end{minipage}%
  \begin{minipage}[c]{\dimexpr\linewidth-4.6cm\relax}\centering
    {\fontsize{28}{33}\selectfont #1}\par\vspace{5pt}
    {\large\color{accent}#2}\par\vspace{8pt}
    {\small #3}\par\vspace{3pt}
    {\small #4}\par
    #5%
  \end{minipage}%
  \begin{minipage}[c]{2.3cm}\raggedleft
    \edef\cv@qr{\noexpand\qrcode[height=2.15cm,level=M]{\cvget{portfolio}}}\cv@qr\par\vspace{2pt}
    {\makebox[2.15cm][c]{\scriptsize\color{sourd}#6}}
  \end{minipage}\par\vspace{4pt}}
\newcommand\cvsep{\hspace{0.6em}{\color{accent!55}\textbar}\hspace{0.6em}}

\makeatother
